% !PRINTED 387
\weberpage{387}

%<<CONT
die Vertauschung von $[1\,4]$ mit $[1\,5]$ in Uebereinstimmung mit der Formel (10). Ebenso leicht erkennt man, dass z.\,B. $[1\,2]$ durch (12) nicht geändert wird.

\begin{center}
\S.\,103.

\gesperrt{Einfachheit der Gruppe des Doppeltangentenproblems.}
\end{center}

Die Darstellung der Gruppe $P$, die wir im vorigen Paragraphen entwickelt haben, liefert uns nun einen ganz einfachen Beweis dafür, dass diese Gruppe \gesperrt{keinen Normaltheiler} hat, also nach unserer früher gebrauchten Ausdrucksweise \gesperrt{einfach} ist, woraus dann folgt, dass die Gruppe nicht durch Adjunction von Irrationalitäten mit kleinerer Gruppe, also beispielsweise nicht durch cyklische Gleichungen erniedrigt werden kann.

Wir stellen $P$ in der Form (4), \S.\,102, dar:
\weq{1}{P = S + \Sigma\, S\, \Pi_{\alpha_1, \alpha_2, \alpha_3, \alpha_4},}
worin $S$ die ganze Gruppe aller Permutationen von acht Ziffern ist. Die Gruppe $S$ hat einen Normaltheiler $S'$ vom Index 2, nämlich die alternirende Gruppe der acht Ziffern, die ihrerseits einfach ist, und jeder Normaltheiler von $S$, der nicht aus der einzigen identischen Permutation besteht, muss die ganze Gruppe $S'$ enthalten. Wir setzen
\weq{2}{S = S' + S'',}
worin $S'' = S'\sigma$ ist, wenn $\sigma$ irgend eine Permutation der zweiten Art, z.\,B. eine Transposition bedeutet.

Wir nehmen jetzt an, es sei $Q$ ein Normaltheiler von $P$, der nicht aus der einzigen identischen Substitution besteht.

Der grösste gemeinschaftliche Theiler von $Q$ und $S$ ist dann ein Normaltheiler von $S$, und muss daher, wenn er nicht die identische Gruppe ist, die Gruppe $S'$ enthalten. Daraus folgt:

1.\ \gesperrt{Wenn $Q$ eine nicht identische Permutation aus $S$ enthält, so enthält $Q$ die ganze Gruppe $S'$.}

Wir beweisen sodann, dass $Q$, wenn es die Gruppe $S'$ enthält, mit $P$ identisch sein muss.

Wenn nämlich $Q$ die ganze Gruppe $S'$ enthält, so enthält
%CONT>>
